% META: {"id": "1SPE-SECDEG-ME-006", "chapitre": "1SPE-SECOND-DEGRE", "type_objet": "methode", "capacites_codes": ["C6"], "methodes": ["M6"], "sources_inspiration": [], "status": "generated"}
\begin{fichemethode}{M6}{Résoudre un problème d'optimisation}

\quandUtiliser{%
  Utiliser cette méthode lorsque l'énoncé demande de « maximiser » ou « minimiser »
  une quantité (aire, volume, bénéfice, périmètre\ldots), de « trouver les dimensions
  qui optimisent », ou de « déterminer la valeur de $x$ pour laquelle la quantité est maximale
  (ou minimale) ».
}

\pasApas{%
  \begin{enumerate}
    \item \textbf{Définir la variable} : je choisis une variable $x$ représentant la quantité qui varie
      et je précise son domaine de définition (les valeurs physiquement admissibles).
    \item \textbf{Exprimer la quantité à optimiser} en fonction de $x$ : j'obtiens une fonction $f(x)$.
    \item \textbf{Identifier le trinôme} : je vérifie que $f$ est un polynôme du second degré
      $f(x) = ax^2 + bx + c$ et je note le signe de $a$.
    \item \textbf{Trouver le sommet} : je complète le carré dans les cas simples
      et je lis $\alpha$ et $\beta$ dans $f(x)=a(x-\alpha)^2+\beta$.
    \item \textbf{Vérifier le domaine et l'objectif} : le sommet donne un minimum
      si $a>0$, un maximum si $a<0$, et seulement s'il appartient au domaine.
      Sur un intervalle réel fermé, je compare les valeurs aux bornes et au sommet
      lorsqu'il est admissible. Une borne exclue ne fournit pas de valeur
      atteinte : sur un intervalle ouvert, l'optimum peut ne pas être atteint.
      Si les valeurs admissibles sont les entiers d'un intervalle borné,
      je compare les valeurs aux entiers admissibles les plus proches du sommet
      et aux deux entiers extrêmes du domaine, selon le minimum ou le maximum recherché.
    \item \textbf{Conclure} en distinguant :
      \begin{itemize}
        \item la \textbf{valeur optimale} retenue (l'aire maximale, le bénéfice maximal\ldots),
        \item la ou les \textbf{valeurs de la variable} qui la réalisent parmi
          les candidats admissibles examinés. Sur un intervalle réel, ce sont
          le sommet ou une borne incluse ; dans le cas discret, ce sont des entiers.
          Si l'optimum n'est pas atteint, je le précise.
      \end{itemize}
  \end{enumerate}
}

\exempleRedige{%
  Dans un repère orthonormé, on considère le triangle de sommets
  $(-4\,;\,0)$, $(4\,;\,0)$ et $(0\,;\,8)$ :
  on inscrit un rectangle symétrique par rapport à l'axe des ordonnées.
  Son côté inférieur est sur l'axe des abscisses et chacun de ses deux
  sommets supérieurs appartient à un côté oblique du triangle.
  Déterminer les dimensions du rectangle d'aire maximale.

  \medskip
  \commentaireMarge{Variable et domaine.}%
  Soit $x$ la demi-largeur du rectangle ($0 < x < 4$).
  La largeur est $2x$. Le sommet supérieur droit appartient à la droite
  d'équation $y=-2x+8$ : la hauteur du rectangle est donc $8-2x$.

  \commentaireMarge{Exprimer l'aire en fonction de $x$.}%
  \[
    A(x) = 2x \times (-2x + 8) = -4x^2 + 16x.
  \]

  \commentaireMarge{Identifier le trinôme : $a = -4 < 0$, donc maximum au sommet.}%
  $A$ est un trinôme avec $a = -4 < 0$, $b = 16$, $c = 0$.

  \commentaireMarge{Compléter le carré et lire le sommet.}%
  \[
    A(x)=-4(x^2-4x)=-4\bigl((x-2)^2-4\bigr)=-4(x-2)^2+16.
  \]
  On lit $\alpha=2$.
  \[
    \beta = A(2) = -4 \times 4 + 16 \times 2 = -16 + 32 = 16.
  \]

  \commentaireMarge{Vérifier que $\alpha = 2 \in \left]0\,;\,4\right[$ : oui.}%
  $\alpha = 2 \in \left]0\,;\,4\right[$ : le sommet est bien dans le domaine.

  \commentaireMarge{Conclure en distinguant variable et valeur optimale.}%
  L'aire est maximale pour $x = 2$, c'est-à-dire une largeur de $2 \times 2 = 4$ et une hauteur
  de $-2 \times 2 + 8 = 4$.

  \medskip
  \textbf{Conclusion :} Le rectangle d'aire maximale a pour dimensions $4 \times 4$,
  et son aire maximale est $A(2) = 16$ unités d'aire.
}

\pieges{%
  \begin{itemize}
    \item \textbf{Piège 1.} Oublier le domaine de validité et l'objectif.
      Un minimum n'est pas un maximum. Sur un intervalle fermé, les bornes
      doivent être examinées ; sur un intervalle ouvert, une borne exclue
      ne peut pas réaliser l'optimum.
    \item \textbf{Piège 2.} Confondre la valeur optimale $\beta = f(\alpha)$ et la variable
      optimale $\alpha$. La question « pour quelle longueur ? » demande $\alpha$ ;
      la question « quelle aire maximale ? » demande $\beta$.
    \item \textbf{Piège 3.} Ne pas définir clairement la variable $x$ avec son domaine dès le début :
      on risque d'obtenir des solutions négatives ou incohérentes avec la situation physique.
  \end{itemize}
}

\verifier{%
  \begin{itemize}
    \item Calculer $A(\alpha - 0{,}5)$ et $A(\alpha + 0{,}5)$ et vérifier qu'ils sont inférieurs
      à $\beta$ (pour un maximum). Ici : $A(1{,}5) = -4 \times 2{,}25 + 24 = 15 < 16$ \checkmark
      et $A(2{,}5) = -4 \times 6{,}25 + 40 = 15 < 16$ \checkmark.
    \item Contrôler que les dimensions trouvées correspondent bien au contexte géométrique
      (largeur et hauteur positives, dans le domaine).
    \item Vérifier les unités de la valeur optimale (aire en unités d'aire, longueur en unités
      de longueur, etc.).
  \end{itemize}
}

\sEntrainer{\refExos{M6}}

\end{fichemethode}

% BEGIN-VERIFY
% from sympy import symbols, expand, Rational
% x = symbols('x', real=True)
% area = 2*x*(8-2*x)
% assert expand(area-(-4*(x-2)**2+16)) == 0
% assert area.subs(x, 2) == 16
% assert area.subs(x, Rational(3, 2)) == area.subs(x, Rational(5, 2)) == 15
% assert 2*(-x)+8 == -2*x+8
% # Le minimum de -x² sur [-1,2] est à la borne2, pas au sommet0.
% assert min(-u*u for u in (-1, 0, 2)) == -4
% # Pour des articles entiers, le sommet 12,5 n'est pas admissible.
% integer_revenue = {n: n*(50-2*n) for n in range(26)}
% assert max(integer_revenue.values()) == 312
% assert [n for n, v in integer_revenue.items() if v == 312] == [12, 13]
% END-VERIFY
